\documentclass[11pt]{article}
\usepackage[utf8]{inputenc}
\usepackage{amsmath}

% simple isotope helper (avoids the isotope package dependency)
\newcommand{\iso}[2]{\textsuperscript{#1}#2}

\begin{document}

\begin{center}
  {\large\textbf{Isotope washout for in-beam PET range verification:}}\\[2pt]
  {\large\textbf{model, evaluation, and implementation}}
\end{center}
\medskip

\noindent
This note specifies the biological washout of the $\beta^+$ emitters and how to
implement it. The model and its window-integrated per-isotope survival factor are
derived first (Secs.~1--4); the implementation recipe for the production chain,
the label-free downstream evaluation actually used here, the per-species selector,
and the one genuine bias route follow (Secs.~5--8). The recommended upstream
implementation is a single Bernoulli keep per recorded decay with probability
$g_i$ (Sec.~5), which is exact for a spatially uniform clearance.

\subsection*{Biological washout model (brain)}

The positron emitters produced by nuclear interactions of the beam
(\iso{15}{O}, \iso{11}{C}, \iso{13}{N}, \dots) do not remain at their
production site: perfusion and metabolism clear the radioactive atoms from the
irradiated volume. Following Mizuno et al.~\cite{mizuno2003}, this biological
washout is described by a sum of three exponential components (fast, medium and
slow), tentatively associated with direct blood flow, the exchange between
tissue and blood vessels, and the trapping of the nuclide by stable molecules,
respectively. The fraction of produced nuclei still present in a voxel a time
$t$ after their creation is
%
\begin{equation}
  W(t) \;=\; \sum_{k=f,m,s} M_k\, e^{-\lambda_k^{\mathrm{bio}}\,t},
  \qquad \sum_{k} M_k = 1,
  \qquad \lambda_k^{\mathrm{bio}} = \frac{\ln 2}{T_k^{\mathrm{bio}}},
  \label{eq:washout}
\end{equation}
%
where $M_k$ is the weight of component $k$ and $T_k^{\mathrm{bio}}$ its
biological half-life. Combining the washout of Eq.~(\ref{eq:washout}) with the
physical decay of the radionuclide $i$ (rate
$\lambda_i^{\mathrm{phys}}=\ln 2 / T_{1/2,i}^{\mathrm{phys}}$), the activity of
a voxel in which $N_0^{\,i}$ nuclei of species $i$ are created at $t=0$ is
%
\begin{equation}
  A_i(t) \;=\; \lambda_i^{\mathrm{phys}}\, N_0^{\,i}
  \sum_{k} M_k\,
  \exp\!\left[-\left(\lambda_i^{\mathrm{phys}}
  + \lambda_k^{\mathrm{bio}}\right) t\right].
  \label{eq:activity}
\end{equation}
%
For a finite irradiation and a given acquisition window, the measured signal is
obtained by convolving the production rate $R_i(t')$ of each species with the
impulse response of Eq.~(\ref{eq:activity}) and integrating over the
acquisition interval $[t_1,t_2]$:
%
\begin{equation}
  S_i \;=\; \int_{t_1}^{t_2}\!\! dt
  \int_{0}^{t}\!\! dt'\; R_i(t')\,
  \lambda_i^{\mathrm{phys}}
  \sum_{k} M_k\,
  \exp\!\left[-\left(\lambda_i^{\mathrm{phys}}
  + \lambda_k^{\mathrm{bio}}\right)(t-t')\right].
  \label{eq:measured}
\end{equation}
%
The washout parameters for brain tissue, as measured in
Ref.~\cite{mizuno2003}, are collected in Table~\ref{tab:washout_brain}.

\begin{table}[ht]
  \centering
  \caption{Three-component biological washout parameters for brain tissue
  (rabbit), from Mizuno et al.~\cite{mizuno2003}. Uncertainties on the last
  digits are given in parentheses. The decay rates
  $\lambda_k^{\mathrm{bio}}=\ln2/T_k^{\mathrm{bio}}$ are listed for convenience
  in both s$^{-1}$ and min$^{-1}$.}
  \label{tab:washout_brain}
  \begin{tabular}{lcccc}
    \hline
    Component & Fraction $M_k$ & $T_k^{\mathrm{bio}}$ (s) &
      $\lambda_k^{\mathrm{bio}}$ (s$^{-1}$) &
      $\lambda_k^{\mathrm{bio}}$ (min$^{-1}$) \\
    \hline
    Fast   & $0.35\,(3)$ & $2.0\,(1.8)$    & $3.5\times10^{-1}$ & $2.1\times10^{1}$ \\
    Medium & $0.30\,(3)$ & $140\,(18)$     & $5.0\times10^{-3}$ & $3.0\times10^{-1}$ \\
    Slow   & $0.35\,(1)$ & $10191\,(2200)$ & $6.8\times10^{-5}$ & $4.1\times10^{-3}$ \\
    \hline
  \end{tabular}
\end{table}

For reference, the physical half-lives of the dominant proton-induced emitters
are $T_{1/2}(\iso{15}{O})=122.2$~s, $T_{1/2}(\iso{13}{N})=597.9$~s and
$T_{1/2}(\iso{11}{C})=1222.8$~s ($20.38$~min).

\subsection*{Evaluation on the frozen source: a per-isotope survival factor}

The productions available to us model physical decay only; each stores, per
species $i$ and depth, the window-integrated signal \emph{without} washout,
which is Eq.~(\ref{eq:measured}) with $W\!\equiv\!1$:
%
\begin{equation}
  S_i^{0} \;=\; \int_{t_1}^{t_2}\!\! dt \int_{0}^{t}\!\! dt'\;
  R_i(t')\,\lambda_i^{\mathrm{phys}}\,
  e^{-\lambda_i^{\mathrm{phys}}(t-t')}.
  \label{eq:phys}
\end{equation}
%
Washout therefore enters the measurement only through the ratio of
Eq.~(\ref{eq:measured}) to Eq.~(\ref{eq:phys}),
%
\begin{equation}
  g_i \;\equiv\; \frac{S_i}{S_i^{0}},
  \label{eq:gdef}
\end{equation}
%
the fraction of physically recorded $i$-decays that also survive clearance.
Both signals carry the same annihilation prefactor $\lambda_i^{\mathrm{phys}}$
and the same production $R_i(t')$ and spatial density, so all of these cancel:
$g_i$ is a pure number per isotope, independent of the absolute normalisation
and --- for a spatially uniform $W$ --- of position.

With a uniform production rate over the irradiation, $R_i(t')=R_i$ for
$t'\in[0,t_{\mathrm{irr}}]$, and an acquisition that opens after the beam
($t_1>t_{\mathrm{irr}}$), the integrals are elementary. Writing
%
\begin{equation}
  \Phi(a) \;\equiv\;
  \frac{\bigl(e^{a\,t_{\mathrm{irr}}}-1\bigr)
        \bigl(e^{-a t_1}-e^{-a t_2}\bigr)}{a^{2}},
  \label{eq:phifun}
\end{equation}
%
one finds $S_i^{0}=R_i\,\lambda_i^{\mathrm{phys}}\,\Phi(\lambda_i^{\mathrm{phys}})$
and $S_i=R_i\,\lambda_i^{\mathrm{phys}}\sum_k M_k\,
\Phi(\lambda_i^{\mathrm{phys}}+\lambda_k^{\mathrm{bio}})$, hence the closed form
%
\begin{equation}
  g_i \;=\; \sum_{k=f,m,s} M_k\,
  \frac{\Phi\!\left(\lambda_i^{\mathrm{phys}}+\lambda_k^{\mathrm{bio}}\right)}
       {\Phi\!\left(\lambda_i^{\mathrm{phys}}\right)} .
  \label{eq:gfactor}
\end{equation}
%
This is a scalar in the Mizuno parameters (Table~\ref{tab:washout_brain}), the
physical rate, and the timing $t_{\mathrm{irr}},t_1,t_2$. The ratio
$\Phi(\lambda+\lambda_k^{\mathrm{bio}})/\Phi(\lambda)$ scales as
$\bigl(\lambda/(\lambda+\lambda_k^{\mathrm{bio}})\bigr)^{2}$ at leading order:
a component that clears far faster than the nuclide decays
($\lambda_k^{\mathrm{bio}}\gg\lambda_i^{\mathrm{phys}}$) contributes nothing,
the atom being gone before it can annihilate. No per-event creation time enters
anywhere --- the age of each decay is integrated out analytically in
Eq.~(\ref{eq:gfactor}).

\subsection*{Exactness of the reweighting}

Applying $g_i$ as a constant thinning of the recorded events reproduces not
only the mean profile but the correct counting noise. Radioactive production is
a Poisson process and clearance acts independently on each nucleus, so by the
thinning theorem the washout-recorded events form a Poisson process whose
intensity is the physical intensity times the survival. Thinning the physical
process by the true age-dependent survival $W(\tau)$ and thinning it by the
constant $g_i$ yield Poisson processes of the \emph{same} intensity $S_i$, and
hence the same distribution --- including the run-to-run spread $\sigma_R$. The
per-event age therefore carries no information the scalar $g_i$ lacks: a plain
Bernoulli keep with probability $g_i$, identical for every event of species
$i$, is exact for the ensemble of independent acquisitions. Three assumptions
underlie this: Poisson production; clearance independent across nuclei (the
single-nucleus picture behind Eq.~(\ref{eq:washout})); and, for a
position-independent $g_i$, spatially uniform $W$.

\subsection*{The measured effect: an isotope-mix reweighting}

Because each $g_i$ multiplies its species' entire depth profile $P_i(z)$ by a
constant, the washed and physical profiles are
%
\begin{equation}
  P^{\mathrm{wo}}(z) = \sum_i g_i\,P_i(z),
  \qquad
  P^{0}(z) = \sum_i P_i(z),
  \label{eq:mix}
\end{equation}
%
so washout distorts no single isotope's shape and acts purely by re-weighting
the mixture. The range endpoint $R_{50}$ moves only insofar as the species
occupy different depths: \iso{15}{O}, with the lowest production threshold,
reaches deepest and owns the distal edge, while \iso{11}{C} lies shallower. The
direction follows from Eq.~(\ref{eq:gfactor}): every recorded nucleus has age
$\ge t_1-t_{\mathrm{irr}}$, far beyond the $2$~s fast scale, so the fast
component is spent for all species and washout roughly halves the signal; the
differential survival comes from the age distribution, and the short-lived
\iso{15}{O} (recorded decays bunched at small age, higher $W$) is suppressed
\emph{less} than the long-lived \iso{11}{C} (decays spread across the window,
more medium and slow clearance). The deep-edge contributor is thus relatively
preserved, and the resulting $R_{50}$ shift partly opposes the proximal walk a
delayed acquisition start produces. Its sign and magnitude follow from
evaluating Eq.~(\ref{eq:gfactor}) for the five species and refitting
$P^{\mathrm{wo}}(z)$ --- a computation on the stored physical-decay columns
alone, with no re-simulation and no per-event data.

\subsection*{Implementation in the production chain}

The exactness result above fixes the recipe, and it is deliberately light: washout
is applied as a \emph{thinning of the already-recorded decay list}, with no
re-simulation of transport and no change to geometry or physics list. In the Monte
Carlo each recorded annihilation carries its emitting species $i$ and its creation
and decay times, which is all that is needed. Two equivalent forms:

\paragraph{Constant per-species thinning (window-integrated) --- recommended.}
Keep each recorded $i$-decay with probability $g_i$ [Eq.~(\ref{eq:gfactor})],
computed once per species from the Mizuno parameters (Table~\ref{tab:washout_brain}),
the physical rate $\lambda_i^{\mathrm{phys}}$, and the timing
$(t_{\mathrm{irr}},t_1,t_2)$. By the thinning theorem this reproduces the exact
washed spatial profile \emph{and} its Poisson noise; no per-event age is used. It
is the simplest complete implementation.

\paragraph{Age-resolved thinning (general).}
Keep each recorded decay with probability $W(\tau)$ [Eq.~(\ref{eq:washout})]
evaluated at its age $\tau=t_{\mathrm{decay}}-t_{\mathrm{creation}}$. For a uniform
$W$ and a fixed window this is statistically identical to the constant-$g_i$ form,
but it needs no recomputation of $g_i$ when the window changes and ---
replacing $W(\tau)\to W(\tau;\mathbf{x})$ with a position-dependent clearance ---
it is the hook for a spatially non-uniform perfusion field (Sec.~8).

Either form is a single Bernoulli keep per event. Because the survival is
position-independent (uniform $W$), it commutes with the detector response: it may
be applied to the truth decay list before photon transport or to the recorded
coincidence list afterwards, with identical result.

\subsection*{Evaluation without the species label}

The reconstruction here is isotope-blind: the detector records only $511$~keV
coincidences, with no species tag, so the per-species $g_i$ cannot be applied
directly. It is recovered by thinning each event with the
isotope-\emph{marginalised} survival
%
\begin{equation}
  w(z_0,t_{\mathrm{d}}) \;=\; \sum_i P(i\,|\,z_0,t_{\mathrm{d}})\; g_i,
  \label{eq:wmarg}
\end{equation}
%
where the posterior that a recorded decay at origin depth $z_0$ and decay time
$t_{\mathrm{d}}$ belongs to species $i$ is
%
\begin{equation}
  P(i\,|\,z_0,t_{\mathrm{d}}) \;=\;
  \frac{P_i(z_0)\,\lambda_i^{\mathrm{phys}}\,
        e^{-\lambda_i^{\mathrm{phys}} t_{\mathrm{d}}}\,\big/\,
        \bigl(1-e^{-\lambda_i^{\mathrm{phys}} t_{\mathrm{meas}}}\bigr)}
       {\sum_j P_j(z_0)\,\lambda_j^{\mathrm{phys}}\,
        e^{-\lambda_j^{\mathrm{phys}} t_{\mathrm{d}}}\,\big/\,
        \bigl(1-e^{-\lambda_j^{\mathrm{phys}} t_{\mathrm{meas}}}\bigr)},
  \label{eq:posterior}
\end{equation}
%
built from the per-species truth profiles $P_i(z)$ and the decay laws, with
$t_{\mathrm{meas}}=t_2-t_1$. Only the origin depth $z_0$ and the decay time
$t_{\mathrm{d}}$ enter --- both recorded per coincidence --- so the marginalised
thinning needs \emph{no} species tag and no new upstream quantity. Averaged over
events it reproduces the mix reweighting Eq.~(\ref{eq:mix}); event by event it is
the label-free surrogate for the constant-$g_i$ keep. With the species label
available upstream, the direct $g_i$ thinning of Sec.~5 is both simpler and exact
and is the recommended path; Eq.~(\ref{eq:posterior}) is what makes the effect
recoverable downstream when the label is absent.

\subsection*{Isolating a single species}

Replacing the marginalised weight Eq.~(\ref{eq:wmarg}) by a single species'
posterior $P(i\,|\,z_0,t_{\mathrm{d}})$ (times a dose factor, no count matching)
selects an ``$i$-like'' sub-sample at its natural abundance. Reconstructing it
yields the per-species range precision $\sigma_R^{(i)}$ --- used, for instance, to
test whether the longer positron range of \iso{15}{O} inflates its $\sigma_R$
relative to \iso{11}{C}. The selection is soft (the sample is enriched, not pure,
so isotope leakage dilutes any contrast into a lower bound).

\subsection*{Spatial non-uniformity: the one genuine bias}

All of the above assumes a spatially uniform clearance $W$, for which washout is a
deterministic, calibratable shift of $R_{50}$ --- it cancels against a per-scanner
reference that itself includes washout. A heterogeneous perfusion field breaks
this: the survival becomes depth-dependent, $g_i\to g_i(z)$, the washed profile is
$\sum_i g_i(z)\,P_i(z)$, and the endpoint shift no longer cancels against a single
constant --- a genuine range bias. Capturing it requires a perfusion/compartment
transport model supplying $W(\tau;\mathbf{x})$ in the age-resolved thinning of
Sec.~5. This is the natural point for an upstream, spatially-resolved treatment,
and is the one part of the effect that the constant-$g_i$ (or marginalised-$w$)
reweighting cannot reproduce.

\begin{thebibliography}{9}

\bibitem{mizuno2003}
H.~Mizuno, T.~Tomitani, M.~Kanazawa, A.~Kitagawa, T.~Kanai, R.~Kohno,
T.~Nishio, et al.,
\newblock Washout measurement of radioisotope implanted by radioactive beams in
the rabbit,
\newblock \emph{Phys. Med. Biol.} \textbf{48}, 2269--2281 (2003).

\end{thebibliography}

\end{document}